\section*{Hauptseminar CS1025 -- Rahmen, Lernziele und Hinweise zur Themenwahl}
\subsection*{Lernziele}
Nach erfolgreicher Teilnahme können Studierende:
\begin{itemize}
  \item aktuelle technische Themen systematisch recherchieren und die Qualität ihrer Quellen bewerten,
  \item einen technischen Gegenstand klar abgrenzen, geeignete Leitfragen formulieren und daraus eine stringente Argumentation ableiten,
  \item wesentliche Sachverhalte zielgruppenorientiert präsentieren und in der Diskussion fundiert vertreten,
  \item Ergebnisse schriftlich nach wissenschaftlichen Standards dokumentieren und kritisch einordnen.
\end{itemize}

\subsection*{Erwartete Leistungen im Seminar}
Das Seminar umfasst typischerweise:
\begin{itemize}
  \item eine Präsentation von ca.\ 30 Minuten,
  \item eine anschließende Diskussion bzw. Fragerunde,
  \item eine schriftliche Ausarbeitung mit ca.\ 4500 Wörtern im Hauptteil,
  \item nachvollziehbare Quellenarbeit, saubere Zitation und eine klare Abgrenzung des Themas.
\end{itemize}

\subsection*{Hinweise zur Themenwahl}
Wähle ein Thema so, dass es fachlich interessant, aber innerhalb des Semesters realistisch bearbeitbar ist.
Hilfreiche Auswahlkriterien sind insbesondere:
\begin{itemize}
  \item ausreichende und belastbare Quellenlage,
  \item ein klar formulierbarer Fokus statt eines zu breiten Überblicks,
  \item ein erkennbarer Bezug zur Ingenieur-Informatik,
  \item ein sinnvoller Ergebnistyp, etwa Vergleich, Einordnung, Systemanalyse oder praxisnahe Bewertung.
\end{itemize}

\noindent\textit{Hinweis:} Die Metadaten jeder Themenkarte (Niveau, Typ, empfohlene Vorkenntnisse, Praxisanteil) dienen als Auswahlhilfe. Nicht jedes Thema erfordert einen praktischen Aufbau; entscheidend sind fachliche Einordnung, saubere Struktur und wissenschaftliche Nachvollziehbarkeit.
